\section*{Capítulo \ref{ch:b3:complex-kinetics} --- Cinética complexa: cadeias, enzimas e oscilações}

\begin{solution}{exo:b3:complex-kinetics:1}
$\ce{Cl2 -> 2 Cl}$: iniciação; $\ce{Cl + H2 -> HCl + H}$ e $\ce{H + Cl2 -> HCl + Cl}$:
propagação; $\ce{2 Cl -> Cl2}$: terminação. Portadores: Cl e H. Reação global (soma das
duas etapas de propagação): $\ce{H2 + Cl2 -> 2 HCl}$.
\end{solution}

\begin{solution}{exo:b3:complex-kinetics:2}
$[\mathrm M]_{1/2} = k_2/k_{-1} = \qty{1.0e-3}{mol/L} = \qty{1.0}{mol/m^3}$; $p = cRT = 1.0 \times 8.314 \times 750 =
\qty{6.2}{kPa}$.
\end{solution}

\begin{solution}{exo:b3:complex-kinetics:3}
$v/V_{\max} = [\mathrm S]/(K_M + [\mathrm S])$: $\frac12$, $\frac34$, $\frac{9}{10}$. Para 90\,\% de $V_{\max}$, é preciso $[\mathrm S] = 9K_M$.
\end{solution}

\begin{solution}{exo:b3:complex-kinetics:4}
$k_{\mathrm{cat}}/K_M = 100/\num{0.80e-3} = \qty{1.25e5}{L.mol^{-1}.s^{-1}}$: cerca de 60\,000 vezes abaixo do limite
de difusão, e próxima do valor típico.
\end{solution}

\begin{solution}{exo:b3:complex-kinetics:5}
Estado estacionário para \ce{CH3CO}: $k_2[\ce{CH3}][\ce{CH3CHO}] = k_3[\ce{CH3CO}]$. Estado estacionário para \ce{CH3}:
$k_1[\ce{CH3CHO}] - k_2[\ce{CH3}][\ce{CH3CHO}] + k_3[\ce{CH3CO}] - 2k_4[\ce{CH3}]^2 = 0$. Somando as duas relações, obtém-se: $k_1[\ce{CH3CHO}] =
2k_4[\ce{CH3}]^2$, logo $[\ce{CH3}] = (k_1/2k_4)^{1/2}[\ce{CH3CHO}]^{1/2}$, e $\dd[\ce{CH4}]/\dd t = k_2[\ce{CH3}][\ce{CH3CHO}] =
k_2(k_1/2k_4)^{1/2}[\ce{CH3CHO}]^{3/2}$.
\end{solution}

\begin{solution}{exo:b3:complex-kinetics:6}
Na metade da conversão, $[\ce{H2}] = [\ce{Br2}] = \frac12$, $[\ce{HBr}] = 1$ (unidades relativas):
$v/v_0 = \frac12 \times (\frac12)^{1/2}/(1 + 0.10 \times 1/0.5) = 0.354/1.2 = 0.29$. Sem inibição, seria
0.354: a inibição divide a velocidade por mais 1.2.
\end{solution}

\begin{solution}{exo:b3:complex-kinetics:7}
Explosão quando $1000p > 6000/p + 10p^2$, isto é, $10p^3 - 1000p^2 + 6000 < 0$, cujas raízes
positivas são 2.48 e \qty{99.9}{kPa}: a mistura explode entre cerca de 2.5 e
\qty{100}{kPa} (primeiro e segundo limites do modelo).
\end{solution}

\begin{solution}{exo:b3:complex-kinetics:8}
$1/v = 1/V_{\max} + (K_M/V_{\max})/[\mathrm S]$: $5.0 = 1/V_{\max} + 2.0K_M/V_{\max}$ e $2.5 = 1/V_{\max} + 0.5K_M/V_{\max}$.
Subtraindo: $K_M/V_{\max} = \qty{1.667}{s.mM.\micro M^{-1}}$, depois $1/V_{\max} = 1.667$: $V_{\max} = \qty{0.60}{\micro M.s^{-1}}$,
$K_M = \qty{1.0}{mM}$.
\end{solution}

\begin{solution}{exo:b3:complex-kinetics:9}
Mesmo intercepto no eixo de $1/v$: competitiva. $K_M^{\mathrm{app}} = 1/0.417 = \qty{2.40}{mM} = 3K_M$, logo
$1 + 50/K_i = 3$ e $K_i = \qty{25}{\micro M}$.
\end{solution}

\begin{solution}{exo:b3:complex-kinetics:10}
$N = \qty{1.01}{mol/L}$: $t_{1/2} = \ln(1.01/0.010 - 1)/(0.50 \times 1.01) = \ln100/0.505 = \qty{9.1}{s}$. A
velocidade $kx(N - x)$ é máxima em $x = N/2$: o mesmo instante, o ponto de inflexão
da curva em S.
\end{solution}

\begin{solution}{exo:b3:complex-kinetics:11}
Traço $B - 2$, determinante 1. $B = 1.5$: $\lambda = -0.25 \pm 0.968\iu$, um foco estável (oscilações
amortecidas em direção ao estado estacionário). $B = 2.5$: $\lambda = 0.25 \pm 0.968\iu$, um foco
instável: as oscilações crescem até atingirem o \omterm{def:b3:complex-kinetics:oscillating}{ciclo limite}.
\end{solution}

\begin{solution}{exo:b3:complex-kinetics:12}
$1/\tau = \num{1.0e8} \times \num{5.40e-5} + \num{1.0e3} = \qty{6.4e3}{s^{-1}}$, $\tau = \qty{156}{\micro s}$. Meça $\tau$ em várias
concentrações totais: $1/\tau$ em função de $[\mathrm A]_e + [\mathrm B]_e$ é uma reta de inclinação $k_1$ e
intercepto $k_{-1}$.
\end{solution}

\begin{solution}{pb:b3:complex-kinetics:1}
\textbf{1.} No início, $[\mathrm S]$ é conhecida e o produto, que poderia inibir ou reagir
de volta, está ausente.
\textbf{2.} Os pontos de $[\mathrm S]$ baixa, em $1/[\mathrm S]$ e $1/v$ grandes: eles fixam a inclinação,
e a inversão amplifica seus erros.
\textbf{3.} Ele pondera cada velocidade tal como foi medida; a reta de Lineweaver--Burk dá
parâmetros enviesados.
\textbf{4.} $K_M$: $0.773$ fica 1.2 incerteza-padrão abaixo de 0.801; $V_{\max}$: $0.491$ fica 2.2
incertezas-padrão abaixo de 0.502. Os dois são baixos, como se espera do viés.
\textbf{5.} $k_{\mathrm{cat}} = 0.502/0.0050 = \qty{100}{s^{-1}}$.
\textbf{6.} $100/\num{0.801e-3} = \qty{1.25e5}{L.mol^{-1}.s^{-1}}$.
\textbf{7.} Cerca de 60\,000 vezes abaixo do limite de difusão; uma \omterm{def:b3:complex-kinetics:enzyme}{enzima} típica.
\textbf{8.} Competitiva: $V_{\max}$ inalterada, $K_M$ aparente crescendo com $[\mathrm I]$.
\textbf{9.} $K_M^{\mathrm{app}} = K_M(1 + [\mathrm I]/K_i) = K_M + (K_M/K_i)[\mathrm I]$.
\textbf{10.} $K_i = 0.799/0.0321 = \qty{24.9}{\micro M}$.
\textbf{11.} $24.9 \pm 4.30 \times 0.9 = 24.9 \pm \qty{3.9}{\micro M}$.
\textbf{12.} A constante de dissociação do complexo enzima--inibidor: a concentração
de inibidor que ocupa metade da \omterm{def:b3:complex-kinetics:enzyme}{enzima} livre.
\textbf{13.} $v_i/v_0 = (K_M + [\mathrm S])/(K_M(1 + [\mathrm I]/K_i) + [\mathrm S])$.
\textbf{14.} $2/3 = 0.67$.
\textbf{15.} $2/12 = 0.17$.
\textbf{16.} $11/21 = 0.52$: em concentração alta, o substrato vence a competição com o
inibidor.
\textbf{17.} Em $[\mathrm S] = K_M$, $v_i/v_0 = 2/(2 + [\mathrm I]/K_i) = 0.10$ dá $[\mathrm I]/K_i = 18$.
\textbf{18.} $18 \times 24.9 = \qty{448}{\micro M}$.
\textbf{19.} A primeira etapa é rápida: cada molécula de \omterm{def:b3:complex-kinetics:enzyme}{enzima} libera depressa um $\mathrm P_1$
e fica presa como \omterm{def:b3:complex-kinetics:enzyme}{acil-enzima}; depois, $\mathrm P_1$ só aparece tão depressa quanto a lenta
segunda etapa libera a \omterm{def:b3:complex-kinetics:enzyme}{enzima}.
\textbf{20.} $1.28/2.0 = 0.64 = (k_2/(k_2 + k_3))^2$, logo $k_2/(k_2 + k_3) = 0.80$.
\textbf{21.} $200/2.0 = \qty{100}{s^{-1}}$, o mesmo valor da Parte I.
\textbf{22.} $k_2 = 0.80 \times 625 = \qty{500}{s^{-1}}$, $k_3 = \qty{125}{s^{-1}}$.
\textbf{23.} $500 \times 125/625 = \qty{100}{s^{-1}}$.
\textbf{24.} \textbf{$[\mathrm I]_{90\,\%} = 18K_i \approx \qty{0.45}{mM}$ em $[\mathrm S] = K_M$.}
\end{solution}
